\section{O Gargalo Fotométrico e a Física da Corrupção do Sinal ($H_2$)}
\label{sec:gargalo_fotometrico}

A confirmação da Hipótese $H_2$, consolidada pelo teste de Friedman na Fase 2, estabelece que a otimização do cenário físico de captura é um preditor de desempenho superior à sofisticação algorítmica isolada. A discussão deste fenómeno exige uma análise sobre como a temperatura de cor da fonte luminosa interage com as propriedades ópticas da fibra de algodão.

A classificação comercial preconizada pela Instrução Normativa nº 24/2016 do MAPA depende criticamente de duas variáveis ópticas: a reflectância ($Rd$) e o grau de amarelamento (+$b$). No \textit{dataset} V1 (Luz Misturada), a coexistência de fontes de 2700 K e 6500 K criou um ambiente de instabilidade espectral. A predominância de comprimentos de onda longos (luz quente) induz um viés fotométrico que mascara a cor real da pluma, ``amarelando'' artificialmente amostras que seriam categorizadas como brancas. Este ruído explica os saltos semânticos detalhados na Secção \ref{sec:analise_qualitativa}, onde a rede confundia classes de grupos distintos (e.g., 44.4 vs. 33.3) devido à incapacidade de distinguir entre o amarelamento intrínseco da fibra e a contaminação cromática da fonte luminosa.

Este achado demonstra que a integridade fotométrica do hardware de captura atua como um fator limitante determinante para a eficácia da inteligência artificial. Modelos de alta capacidade paramétrica têm seu potencial de generalização reduzido quando operam sobre sinais de entrada degradados. Portanto, a padronização fotométrica não é apenas uma recomendação normativa, mas uma exigência técnica para a viabilidade de sistemas de \textit{Edge AI} no processamento de \textit{commodities}.